\begin{table}[htbp]
\centering
\caption{Frequent-Loser Price Coefficient by Procuring-Unit Size and Procedure}
\label{tab:regime_oversight}
\begin{adjustbox}{max width=\textwidth}
\begin{threeparttable}
\small
\begin{tabular}{lccc}
\toprule
Sub-sample & Coefficient & SE & N \\
\midrule
\textit{By procuring-unit size quartile:} & & & \\
\quad Q1 (smallest buyers) & 0.2138 & (0.1670) & 7,503 \\
\quad Q2 & 0.0111 & (0.0513) & 87,571 \\
\quad Q3 & 0.0660** & (0.0304) & 302,258 \\
\quad Q4 (largest buyers) & 0.0165 & (0.0160) & 1,257,069 \\
\midrule
\textit{By procedure type:} & & & \\
\quad Preg\~{a}o (reverse electronic auction) & 0.0933*** & (0.0255) & 546,549 \\
\quad Convite (sealed-bid invitation) & 0.0382** & (0.0186) & 1,107,852 \\
\bottomrule
\end{tabular}
\begin{tablenotes}
\small
\item \textit{Notes:} A coefficient that shrinks across larger buyers is
consistent with thinner institutional capacity at smaller buyers
amplifying the participation footprint the construct flags; we do not
identify oversight as the operative channel
(\S\ref{sec:results_oversight}). Item and year FE (PBU FE for procedure
split). SE clustered at item level.
*** \textit{p}$<$0.01, ** \textit{p}$<$0.05, * \textit{p}$<$0.1.
\end{tablenotes}
\end{threeparttable}
\end{adjustbox}
\end{table}
